% =====================================================================
%  AI4Good workshop variant.
%  Defines everything main.tex pulls in: title, venue string, abstract,
%  introduction, discussion. Nothing else in the paper changes.
% =====================================================================

\newcommand{\wkTitle}{Portable Causal Fairness Across \\
Synthetic Data Generator Families}

\newcommand{\wkVenue}{AI for Good}

% ---------------------------------------------------------------- abstract
\newcommand{\wkAbstract}{%
Institutions that release synthetic data in place of sensitive records, such as
a hospital, a lending regulator or a statistical agency, choose the process
that generates the table, and can shape that process so unfair pathways are
absent. DECAF made this concrete for one GAN: factor the generator along a
causal graph, sever the edges from a protected attribute to an outcome, and
three fairness definitions become three sets of edge cuts. But DECAF is a
single GAN, carrying no privacy guarantee and no competitive data-quality
story, and the mechanism has never been shown to work in any other
architecture. We ask whether it belongs to DECAF or to causal factorisation
itself, porting all three definitions to nine generators from three unrelated
families (marginals-based, GAN-based and diffusion-based, each with
differentially private variants) across 2{,}520 matched-pair runs on the Adult
and COMPAS datasets. The definitions transfer everywhere: cutting edges
reduces downstream disparity in every family, dataset and privacy level. The
utility cost is small (under two accuracy points for the cuts worth making in
five of six family--dataset cells), and the privacy cost is nil in the
marginals family, because a severed edge is never measured under the privacy
budget. In this setting privacy and fairness turn out to behave close to
independently rather than adversarially. Causal fairness is now available
inside a modern differentially private synthesizer and inside a diffusion
backbone, not only inside one GAN.%
}

% ------------------------------------------------------------------- intro
\newcommand{\wkIntro}{%
Institutions holding sensitive records increasingly publish a synthetic
substitute rather than the records themselves: a hospital, a lending regulator
or a statistical agency fits a generative model and releases a table drawn from
it. That substitution is usually justified on privacy grounds, but it carries a
second, less-exploited consequence. Whoever generates the data chooses the
generating process, and a process can be edited. A disparity that exists in the
world because of a pathway running from a protected attribute, such as national
origin, to an outcome, such as whether a loan is approved, need not exist in
the released table.

The cleanest existing realisation of that lever is
DECAF~\citep{vanbreugel2021decaf}, which factors a generator along a causal
graph so that severing an edge at sampling time makes the generator
structurally unable to use the pathway it named. But DECAF is a single
non-private WGAN-GP~\citep{gulrajani2017wgangp}, whose released tables trail
modern differentially private marginals-based synthesizers and modern
diffusion models on the fidelity and utility benchmarks that decide what
institutions actually deploy. Whether the fairness mechanism is confined to
that one architecture, or is a property of causal factorisation itself, is
therefore a practical question and not an academic one: it decides whether an
institution deploying synthetic data has the intervention available at all.

We ask two questions of that mechanism. Is edge surgery a property of DECAF,
or of the causal factorisation any such generator uses? And can it be attached
to a generator whose graph is not a supplied causal DAG but the undirected
marginal-dependency graph a differentially private synthesizer builds from
noisy counts? The stakes are practical: if the mechanism that would make a
released table fair is confined to a generator no institution would actually
deploy, it might as well not exist.%
}

% -------------------------------------------------------------- discussion
\newcommand{\wkDiscussion}{%
\textbf{Edge surgery belongs to causal factorisation, not to DECAF.} DECAF
conjectured that its mechanism would extend beyond its own GAN but tested no
other generator. It extends: across nine generators from three unrelated
families, both datasets and every privacy budget, severing edges reduces
downstream disparity, and FTU, DemPar and CF compile to the same edge cuts on
a diffusion backbone and on a differentially private marginals synthesizer as
they do on the original GAN. An institution can now choose a synthesizer for
its fidelity, privacy or maturity and keep the fairness lever, instead of being
pushed onto a non-private GAN to get it.

\textbf{Privacy and fairness barely interact here, and in the marginals family
the fairness cut costs no privacy at all.} The textbook expectation is a
trade-off, with tighter privacy blunting any attempt to close the parity gap.
We see the opposite: the cut removes at least as much disparity under the tight
$\varepsilon{=}1$ budget as under a permissive one. And in a marginals-based
synthesizer the cut is free in the privacy accounting, because an edge the
definition forbids is simply never measured under the budget, a property
DECAF's GAN cannot offer since it pays for every column it draws. For an
institution already publishing under a privacy guarantee, fairness here is
close to a free addition rather than a competing demand.

\textbf{How much the cut removes is set by the graph, not the generator.} The
definition is identical everywhere; its effect size is not, for a simple
reason. A hand-supplied causal DAG routes a protected attribute to the outcome
through many edges, leaving a large disparity to sever; a private marginals
graph is sparse by construction and often carries that dependence faintly or
not at all. On Adult, FTU on MST is a literal no-op in all $45$ matched pairs,
because the private structure search never buys the direct
\texttt{sex}$\rightarrow$\texttt{income} edge in the first place: the
constraint is satisfied before it is imposed. So the question an auditor should
ask of a release is not which generator produced it, but which graph it was
sampled along and which attributes were nominated admissible.

\textbf{Only the in-generator cut ships with the table.} Where a private
structure encodes little to sever, as on COMPAS, downstream mitigators remain
the stronger tool. But a public table is used by consumers its publisher will
never meet, and of the four places to intervene the in-generator cut is the
only one that travels with the release and cannot be dropped by whoever trains
the next model on it. The two are complements: the publisher should apply the
cut wherever the graph gives it something to sever.

\textbf{Conclusion.} Synthetic data is increasingly the vehicle through which
institutions share what they know for the public good, as when a hospital, a
regulator or a statistical agency releases a table so that others can build on
it. Our results show that the fairness of such a release need not be left to
whoever eventually trains a model on it, nor traded against the privacy and
fidelity these institutions rightly demand: a single causal edit, portable
across the synthesizers they actually deploy, can remove an unfair pathway
before the data ever leaves the building. Making that lever available wherever
data is generated is a small but concrete step toward synthetic releases that
are fair by construction.%
}
